\subsection{幺星代数中的函数演算}

在本小节中，A 是一个幺星代数，x 是 A 的一个正规元。

设 B 是 A 的由 x 生成的幺星子代数；它是交换的，且包含在 \(\{x, x^{*}\}\) 的双交换子中，从而包含在 x 的双交换子中 (I, p.~106 命题 6 的推论)。Gelfand 变换 \(\mathcal{G}_{\mathrm{B}}: \mathrm{B} \to \mathcal{C}(\mathrm{X}(\mathrm{B}))\) 是星代数的一个同构 (I, p.~108, 定理 1)。我们有 \(\mathrm{Sp}_{\mathrm{B}}(x) = \mathrm{Sp}_{\mathrm{A}}(x)\) (I, p.~106, 命题 5)。

引理 10. --- 映射 \(\mathrm{ev}_x\colon \chi \mapsto \chi (x)\) 诱导出从 \(\mathsf{X}(\mathbf{B})\) 到 \(\mathrm{Sp}_{\mathrm{A}}(x)\) 上的一个同胚。

\(x \mapsto \chi(x)\) 是从 \(\mathsf{X}(\mathsf{B})\) 到 \(\mathbf{C}\) 中的一个连续映射，依 I, p.~37 的命题 6，它的像等于 \(\mathrm{Sp}_{\mathsf{B}}(x)\)，从而等于 \(\mathrm{Sp}_{\mathsf{A}}(x)\)。由于 B 的特征标都是 Hermite 的 (I, p.~107, 引理 9)，在 \(x\) 上取值相同的 B 的特征标在 \(x^{*}\) 上也取相同的值，因而在由 \(x\) 生成的幺星子代数 B 上相等。这就证明映射 \(\mathrm{ev}_x\) 是单射。由于 \(\mathsf{X}(\mathsf{B})\) 是紧的且 \(\mathbf{C}\) 是分离的，它诱导出从 \(\mathsf{X}(\mathsf{B})\) 到其像上的一个同胚，由此得引理。

由该引理我们导出星代数的一个同构 \(\varphi_x\colon \mathcal{C}(\mathrm{Sp}_{\mathrm{A}}(x))\to \mathcal{C}(\mathrm{X(B)})\)。它把函数 \(f\in \mathcal{C}(\mathrm{Sp}_{\mathrm{A}}(x))\) 映为函数 \(\chi \mapsto f(\chi (x))\)。映射 \(\mathcal{G}_{\mathrm{B}}^{-1}\circ \varphi_{x}\) 是从星代数 \(\mathcal{C}(\mathrm{Sp}_{\mathrm{A}}(x))\) 到 B 上的一个同构。

定义 5. --- 对合态射 \(f \mapsto (\mathcal{G}_{\mathrm{B}}^{-1} \circ \varphi_x)(f)\)（由 \(\mathcal{C}(\mathrm{Sp}_{\mathrm{A}}(x))\) 到 A）称为 \(x\) 在 A 中的连续函数演算映射。我们把它记作 \(f \mapsto f(x)\)。

注记. --- 映射 \(f \mapsto f(x)\) 是等距的；它的像是由 x 生成的幺星子代数 B，它包含在 x 的双交换子中。

若 \(f\) 在 \(\mathrm{Sp}_{\mathrm{A}}(x)\) 上的限制是形如 \(z \mapsto \mathrm{P}(z, \overline{z})\) 的函数（其中 \(\mathrm{P} \in \mathbf{C}[\mathrm{X}, \mathrm{Y}]\) 是一个多项式），则按通常的代数意义有 \(f(x) = \mathrm{P}(x, x^{*})\)。

例. --- 1) 设存在 \(\lambda \in \mathbf{C}\) 使 \(\mathrm{Sp}_{\mathrm{A}}(x)\) 缩减为 \(\lambda\)。那么我们有 \(x = \lambda \cdot 1\)。事实上，\(\mathrm{Sp}_{\mathrm{A}}(x)\) 上的恒等函数等于 \(\lambda\)，因而它在函数演算映射下的像，即 \(x\)，等于 \(\lambda \cdot 1\)。

\begin{enumerate}
\def\labelenumi{\arabic{enumi})}
\setcounter{enumi}{1}
\item
  为使 \(x\) 是 Hermite 的，必须且只需 \(\mathrm{Sp}_{\mathrm{A}}(x)\) 包含在 \(\mathbf{R}\) 中。事实上，设 \(f\) 是 \(\mathrm{Sp}_{\mathrm{A}}(x)\) 上由 \(f(z) = z - \overline{z}\) 给定的连续映射。那么 \(x\) 是 Hermite 的当且仅当 \(f(x) = 0\)，即当且仅当 \(f\) 为零，亦即当且仅当 \(\mathrm{Sp}_{\mathrm{A}}(x)\) 包含在 \(\mathbf{R}\) 中。
\item
  为使 \(x\) 是酉元，必须且只需它的谱包含在 \(\mathbf{C}\) 的单位圆中。事实上，设 \(f \in \mathcal{C}(\mathrm{Sp}_{\mathrm{A}}(x))\) 是由 \(f(z) = z\overline{z} - 1\) 定义的函数；元素 \(x\) 是酉元当且仅当 \(f(x) = 0\)，即当且仅当 \(f\) 为零。
\end{enumerate}

命题 7. — 映射 \(f \mapsto f(x)\) 是从 \(\mathcal{C}(\mathrm{Sp}_{\mathrm{A}}(x))\) 到 A 中的唯一的幺对合代数态射，使得恒等映射 \(z\)（在 \(\mathrm{Sp}_{\mathrm{A}}(x)\) 上）以 \(x\) 为像。

事实上，由 \(\mathcal{C}(\mathrm{Sp}_{\mathrm{A}}(x))\) 的元素 \(z\) 与 \(\overline{z}\) 生成的 \(\mathcal{C}(\mathrm{Sp}_{\mathrm{A}}(x))\) 的幺子代数在 \(\mathcal{C}(\mathrm{Sp}_{\mathrm{A}}(x))\) 中稠密 (TG, X, p.~40, 推论 1)。由于从 \(\mathcal{C}(\mathrm{Sp}_{\mathrm{A}}(x))\) 到 A 的每个对合代数态射都是连续的 (I, p.~104, 命题 2)，至多存在一个从 \(\mathcal{C}(\mathrm{Sp}_{\mathrm{A}}(x))\) 到 A 的对合代数态射，它把 \(z\) 映为 \(x\)。

下面的推论表明：当 \(f\) 是 \(\mathrm{Sp}_{\mathrm{A}}(x)\) 的一个邻域中全纯函数的限制时，\(f(x)\) 的这一定义与 I, p.~74 第 9 小节中单变量全纯函数演算的定义一致。

推论 1. --- 设 \(f \in \mathcal{O}(\mathrm{Sp}_{\mathrm{A}}(x))\) 是 \(\mathrm{Sp}_{\mathrm{A}}(x)\) 的一个邻域中全纯函数的一个芽，\(\widetilde{f} \in \mathcal{C}(\mathrm{Sp}_{\mathrm{A}}(x))\) 是 \(\mathrm{Sp}_{\mathrm{A}}(x)\) 上的与 \(f\) 相伴随的连续函数。我们有 \(\widetilde{f}(x) = f(x)\)，其中 \(f(x)\) 是由全纯函数演算给出的 A 的元素。

事实上，映射 \(f \mapsto \widetilde{f}(x)\) 是从 \(\mathcal{O}(\mathrm{Sp}_{\mathrm{A}}(x))\) 到 A 的一个连续幺态射，它把 \(\mathrm{Sp}_{\mathrm{A}}(x)\) 的一个邻域中恒等函数的芽映为 \(x\)。于是结果由 I, p.~74 的定理 5 得出。

推论 2. --- 设 \(f \in \mathcal{C}(\mathrm{Sp}_{\mathrm{A}}(x))\)。

\begin{enumerate}
\def\labelenumi{\alph{enumi})}
\tightlist
\item
  我们有
\end{enumerate}

\[
\operatorname{Sp} _ {\mathrm{A}} (f (x)) = f (\operatorname{Sp} _ {\mathrm{A}} (x));
\]

\begin{enumerate}
\def\labelenumi{\alph{enumi})}
\setcounter{enumi}{1}
\tightlist
\item
  对每个 \(g \in \mathcal{C}(\mathrm{Sp}_{\mathrm{A}}(f(x)))\)，有 \((g \circ f)(x) = g(f(x))\)。
\end{enumerate}

由于 \(f(x)\) 属于 A 的全子代数 B，我们有 \(\mathrm{Sp}_{\mathrm{A}}(f(x)) = \mathrm{Sp}_{\mathrm{B}}(f(x))\) (I, p.~106, 命题 5)。同构 \(f \mapsto f(x)\)（从 \(\mathcal{C}(\mathrm{Sp}_{\mathrm{A}}(x))\) 到 B）保持谱；因此我们有 \(\mathrm{Sp}_{\mathrm{B}}(f(x)) = \mathrm{Sp}_{\mathcal{C}(\mathrm{Sp}_{\mathrm{A}}(x))}(f) = f(\mathrm{Sp}_{\mathrm{A}}(x))\) (I, p.~17, 例 3)。这就证明了断言 \(a\))。

映射 \(g \mapsto (g \circ f)(x)\) 是从 \(\mathcal{C}(\mathrm{Sp}_{\mathrm{A}}(f(x)))\) 到 A 的一个幺对合代数态射，它把 \(\mathrm{Sp}_{\mathrm{A}}(f(x))\) 的恒等函数映为 \(f(x)\)。于是依命题 7，\((g \circ f)(x) = g(f(x))\) 对一切 \(g \in \mathcal{C}(\mathrm{Sp}_{\mathrm{A}}(f(x)))\) 成立。

例 4. --- 设 X 是一个局部紧空间，A = \(\mathcal{C}_{b}(X)\) 是 X 上有界连续函数组成的交换幺星代数 (I, p.~102, 例 2)。设 \(g \in A\)；它的谱 S 是 g 的值集 \(g(\mathrm{X})\) 在 C 中的闭包 (I, p.~17, 例 3)。于是 g 的函数演算映射就是映射 \(f \mapsto f \circ g\)（对 \(f \in \mathcal{C}(\mathrm{S})\)）。事实上，这个映射是星代数的一个幺态射，且使 S 的恒等映射以 g 为像。

在 X 是紧的情形，我们有 \(A = \mathcal{C}(X)\) 且 \(S = g(X)\)。

设 \(\pi\colon\mathrm{A}\to\mathrm{A}^{\prime}\) 是幺星代数的一个幺态射。元素 \(\pi(x)\) 在 \(\mathrm{A}^{\prime}\) 中是正规的，且它相对于 \(\mathrm{A}^{\prime}\) 的谱含于 \(\mathrm{Sp}_{\mathrm{A}}(x)\) 中。于是我们有：

命题 8. --- 设 \(f \in \mathcal{C}(\mathrm{Sp}_{\mathrm{A}}(x))\)。仍用 \(f\) 表示 \(f\) 在 \(\mathrm{Sp}_{\mathrm{A}'}(\pi(x))\) 上的限制，我们有等式 \(\pi(f(x)) = f(\pi(x))\)。特别地，对每个 \(\chi \in \mathsf{X}(\mathrm{A})\)，有 \(\chi(f(x)) = f(\chi(x))\)。

设 \(z\) 是 \(\mathrm{Sp}_{\mathrm{A}}(x)\) 的恒等映射。由 \(f \mapsto \pi(f(x))\) 与 \(f \mapsto f(\pi(x))\) 定义的映射都是从 \(\mathcal{C}(\mathrm{Sp}_{\mathrm{A}}(x))\) 到 B 的连续幺对合代数态射，且都把 \(z\) 映为 \(\pi(x)\)。因此这两个态射在 \(\mathcal{C}(\mathrm{Sp}_{\mathrm{A}}(x))\) 的由 \(z\) 生成的幺对合子代数上一致。由于后者在 \(\mathcal{C}(\mathrm{Sp}_{\mathrm{A}}(x))\) 中稠密 (TG, X, p.~40, 推论 1)，这两个态射相等。

推论. --- 设 A 是交换的。对每个 \(f \in \mathcal{C}(\mathrm{Sp}_{\mathrm{A}}(x))\)，有 \(\mathcal{G}_{\mathrm{A}}(f(x)) = f \circ \mathcal{G}_{\mathrm{A}}(x)\)。

只需对 Gelfand 变换 \(\mathcal{G}_{\mathrm{A}}\)（从 A 到 \(\mathcal{C}(\mathsf{X}(\mathrm{A}))\)）应用命题 8，并注意（见上面的例）\(f(\mathcal{G}_{\mathrm{A}}(x)) = f \circ \mathcal{G}_{\mathrm{A}}(x)\)。

